\chapter{开放系统}
\label{ch:open}

\section{约化动力学与量子信道}

上一章用干涉条纹的位置读出相位，前提是两路振幅还能相干叠加。实际把原子制备成叠加态，
延长等待时间，还会看到 Ramsey 条纹的对比度下降。若把同种原子制备到激发态，寿命测量记录的激发态比例也会下降，
两种衰减的速率可以不同。若要用同一模型解释两组读数，需要确定光子离开后，仅测量原子还能获得哪些信息。

从封闭总系统出发，可以保留概率守恒而只观测其中一部分。
开放系统(open quantum system)与外部自由度交换能量或信息。
把关心的自由度记为系统$\mathrm{S}$，未逐项观测的自由度记为环境(environment) $\mathrm{E}$。
系统的状态是联合密度算符对环境的部分迹，
$\op{\rho}_{\mathrm{S}} = \tr_{\mathrm{E}}\op{\rho}_{\mathrm{SE}}$。
例如发光原子是$\mathrm{S}$，逸出到各个方向的辐射模式组成$\mathrm{E}$。

\begin{example}{自发发射的光子计数与原子态}{qm-amplitude-damping}
    设辐射场最初处于真空态 $\ket{0}_{\mathrm{E}}$，发射一个光子后的归一化场态记为 $\ket{1}_{\mathrm{E}}$。
    按原子、辐射场的顺序写乘积基，例如 $\ket{\excitedstate,0} = \ket{\excitedstate}_{\mathrm{S}} \otimes \ket{0}_{\mathrm{E}}$ 表示原子在激发态而场中没有光子。
    在一次演化区间内，以$p \in [0,1]$表示激发态发射的概率。
    联合演化在相关初态子空间上为
    \begin{subequations}
        \begin{alignat}{1}
            \ket{\groundstate,0} &\longmapsto \ket{\groundstate,0},\\
            \ket{\excitedstate,0} &\longmapsto
                \sqrt{1 - p}\ket{\excitedstate,0} + \sqrt p\ket{\groundstate,1}.
        \end{alignat}
    \end{subequations}
    设原子的初态为 $a\ket{\groundstate} + b\ket{\excitedstate}$，其中 $\abs{a}^2 + \abs{b}^2 = 1$。
    求记录零个或一个光子的概率，再求不读取光子记录时任意原子初态的演化。

    \par\smallskip\textbf{解}\quad 两个输出仍正交归一，因此可扩充为整个联合空间上的幺正变换。根据线性叠加，等候结束后的联合态为
    \begin{equation}
        \bigp{a\ket{\groundstate}_{\mathrm{S}} + \sqrt{1 - p}\,b\ket{\excitedstate}_{\mathrm{S}}} \otimes \ket{0}_{\mathrm{E}}
            + \sqrt p\,b\ket{\groundstate}_{\mathrm{S}} \otimes \ket{1}_{\mathrm{E}} .
    \end{equation}
    光子计数器理想时，这两个正交的环境记录可以区分。读到零个光子的概率是第一项的模方
    $\abs{a}^2 + (1 - p)\abs{b}^2$，读到一个的概率是第二项的模方 $p\abs{b}^2$。
    例如初态为等幅叠加、$p = 1/2$，四分之一的实验会记录一个光子，原子随即确定在基态。
    其余四分之三记录为零，其中原子仍在激发态的条件概率是 $(1/4)/(3/4) = 1/3$。

    为了将计算推广到任意原子初态，可分别考察零光子和一个光子记录下的原子态。它们对初态的线性依赖分别由下列两个算符给出，
    \begin{equation}
        \op{K}_0 = \ketbra{\groundstate}{\groundstate} + \sqrt{1 - p}\ketbra{\excitedstate}{\excitedstate},
        \qquad
        \op{K}_1 = \sqrt p\ketbra{\groundstate}{\excitedstate}.
    \end{equation}
    若只测量原子，各条光子记录所对应的条件态应按记录概率加权。
    条件态的归一化分母与权重相消。以 $(\ket{\groundstate},\ket{\excitedstate})$ 为有序基，把两种记录的贡献相加，
    \begin{equation}
        \begin{alignedat}{1}
            \op{\rho}' &= \op K_0\op\rho\op K_0\adjoint + \op K_1\op\rho\op K_1\adjoint\\
            &= \begin{pmatrix}
                \rho_{\groundstate\groundstate} + p\rho_{\excitedstate\excitedstate}&\sqrt{1 - p}\rho_{\groundstate\excitedstate}\\
                \sqrt{1 - p}\rho_{\excitedstate\groundstate}&(1 - p)\rho_{\excitedstate\excitedstate}
            \end{pmatrix}.
        \end{alignedat}
    \end{equation}
    激发态布居减少因子$1 - p$，相干只减少因子$\sqrt{1 - p}$，
    因为相干连接“尚未发射”的振幅与基态振幅。
    初态为$\ket{\excitedstate}$时，约化态为
    $(1 - p)\ketbra{\excitedstate}{\excitedstate} + p\ketbra{\groundstate}{\groundstate}$；
    混合来自系统与光子的关联，联合态仍是纯态。
\end{example}

上述两种记录可以推广到任意环境基，即使实际实验只观测原子也可以求和。
封闭联合系统满足$\op{\rho}_{\mathrm{SE}}(t) = \op{U}\op{\rho}_{\mathrm{SE}}(0)\op{U}\adjoint$。
若初始态可分为$\op{\rho}_{\mathrm{S}} \otimes \op{\rho}_{\mathrm{E}}$且环境处于纯态$\ket{e_0}$，选环境基$\ket{\alpha}$，定义
\begin{subequations}
    \begin{align}
        \op{K}_\alpha &= (\one_{\mathrm{S}} \otimes \bra{\alpha})\op{U}(\one_{\mathrm{S}} \otimes \ket{e_0}),\\
        \op{\rho}_{\mathrm{S}}(t) &= \sum_\alpha \op{K}_\alpha\op{\rho}_{\mathrm{S}}\op{K}_\alpha\adjoint,\\
        \sum_\alpha \op{K}_\alpha\adjoint \op{K}_\alpha &= \one_{\mathrm{S}}.
    \end{align}
\end{subequations}

$\op{K}_\alpha$ 称为 Kraus 算符，作用在系统空间上。右端先把系统态与 $\ket{e_0}$ 配成初态，
联合演化后，左端只取环境态 $\ket{\alpha}$ 的分量，留下一个系统态。于是可简记 $\op{K}_\alpha = \bra{\alpha}\op{U}\ket{e_0}$。
对初始纯系统态$\ket{\psi}$，联合演化可展开为
\begin{equation}
    \op{U}(\ket{\psi} \otimes \ket{e_0})
    = \sum_\alpha(\op{K}_\alpha\ket{\psi}) \otimes \ket{\alpha}.
\end{equation}
求外积并取环境的迹，因环境基正交，具有不同 $\alpha$ 的交叉项消失，余下各项给出上式中的约化态。
由线性，该式对一般$\op{\rho}_{\mathrm{S}}$同样成立。
再由环境完备性和$\op{U}\adjoint \op{U} = \one$，
\begin{equation}
    \sum_\alpha \op{K}_\alpha\adjoint \op{K}_\alpha = (\one_{\mathrm{S}} \otimes \bra{e_0})
        \op{U}\adjoint
        \Bigp{\one_{\mathrm{S}} \otimes \sum_\alpha\ketbra{\alpha}{\alpha}}
        \op{U}(\one_{\mathrm{S}} \otimes \ket{e_0}) = \one_{\mathrm{S}}.
\end{equation}

若 $\op{\rho}_{\mathrm{E}}$ 是混合态，取其正交本征态 $\ket{e_j}$，以本征值 $w_j$ 为各初始环境态的权重。
将权重吸收到 $\op{K}_{\alpha j} = \sqrt{w_j}\bra{\alpha}\op{U}\ket{e_j}$ 中，仍得到上述形式。
这个作用于任意 $\op{\rho}_{\mathrm{S}}$ 的映射来自固定环境态的乘积制备。若初态包含系统与环境的关联，预测演化还须指定这些关联。

\begin{definition}{完全正迹保持映射}{qm-cptp}
    线性映射$\mathcal E$对所有输入态均满足$\tr\mathcal E(\op{\rho}) = \tr\op{\rho}$，称为迹保持(trace-preserving)。
    若对任意辅助系统 $\mathrm{A}$，$\mathcal E \otimes \id_{\mathrm{A}}$都把正算符映到正算符，
    称为完全正(completely positive)。同时满足这两条的映射称为量子信道(quantum channel)，
    英文简写为 CPTP (completely positive and trace-preserving)。
\end{definition}
例如对一对纠缠粒子的其中一个施加信道，另一个置于装置之外，所有联合概率仍须为非负。
完全正性保证这种只在一侧操作的实验也有一致的概率解释。

\begin{theorem}{Kraus算符和的完全正性与保迹条件}{qm-kraus-trace}
    有限维中，给定一组算符 $\op{K}_\alpha$，令
    \begin{subequations}\label{eq:qm-kraus-channel-criterion}
        \begin{align}
            \mathcal E(\op{\rho}) &= \sum_\alpha \op{K}_\alpha\op{\rho} \op{K}_\alpha\adjoint,\\
            \op{Q} &= \sum_\alpha \op{K}_\alpha\adjoint \op{K}_\alpha .
        \end{align}
    \end{subequations}
    映射 $\mathcal E$ 完全正。它是量子信道的充要条件为 $\op{Q} = \one$。
\end{theorem}
\begin{proof}
    由迹的循环性质，$\tr\mathcal E(\op{\rho}) = \tr(\op{\rho} \op{Q})$。
    当 $\op{Q} = \one$ 时便有迹保持。反之，在每个归一化纯态上测试迹保持，
    就得到自伴算符 $\op{Q}$ 在每个纯态中的期望值均为 1，故 $\op{Q} = \one$。
    对联合正算符$\op{R} \geq 0$及任意$\ket{v}$，
    \begin{subequations}
        \begin{align}
            \mel{v}{(\mathcal E \otimes \id_{\mathrm{A}})(\op{R})}{v}
            &= \sum_\alpha
            \mel{w_\alpha}{\op{R}}{w_\alpha} \geq 0,\\
            \ket{w_\alpha} &= (\op{K}_\alpha\adjoint \otimes \one_{\mathrm{A}})\ket{v}.
        \end{align}
    \end{subequations}
    因此加入任意辅助空间后仍保持正性。
\end{proof}
取部分迹后，环境携带的路径或发射信息无法仅从原子读出。由系统与环境关联引起的局部相干衰减称为退相干(decoherence)，
在干涉实验中表现为可见度下降。对光子记录求平均得到约化态，保留某个记录则得到对应的条件态。

\section{Lindblad生成元与无记忆演化}

上节的信道给出原子等待一段时间后的状态。现在把等待分成两段。若第一段发出的光会返回原子，
第二段的演化就可能受早先的发射影响。若环境记忆可以忽略且外部条件保持恒定，只需给定当前的原子态，便可用同一演化规则预测它在下一时段的状态。
在这个条件下，一条微分方程足以描述整条衰减曲线。

\begin{definition}{时间齐次动力学半群与生成元}{qm-semigroup}
    用 $\mathcal E_t$ 表示经过时间 $t \geq 0$ 的量子信道。若
    $\mathcal E_0 = \id$ 且 $\mathcal E_{t+s} = \mathcal E_t\mathcal E_s$，这族信道称为时间齐次动力学半群(semigroup)。
    初始变化率的极限若存在，便定义生成元(generator) $\mathcal L$。由合成律，将它作用于时刻 $t$ 的态就得到主方程(master equation)，
    \begin{subequations}
        \begin{align}
            \mathcal L(\op{\rho}) &= \lim_{h\to0^+}\frac{\mathcal E_h(\op{\rho}) - \op{\rho}}{h},\\
            \dv{\op{\rho}(t)}{t} &= \mathcal L\bigp{\op{\rho}(t)}.
        \end{align}
    \end{subequations}
\end{definition}

以\exmref{exm:qm-amplitude-damping}为例，设$p = p(t)$。
连续两段演化使未发射比例相乘，半群条件要求
\begin{equation}
    1 - p(t + s) = [1 - p(t)][1 - p(s)].
\end{equation}
若$p(0) = 0$且随时间连续，该乘法关系给出
$1 - p(t) = \e^{-\Gamma t}$，其中$\Gamma \geq 0$。
具体说，设激发态在一个短的时间段 $\Delta t$ 内有 $\Gamma\Delta t + o(\Delta t)$ 的概率发射。
到 $t$ 时仍激发的比例记为 $s(t)$，则 $s(t + \Delta t) = s(t)[1 - \Gamma\Delta t + o(\Delta t)]$。
移项除以 $\Delta t$ 就得到 $\dot s = -\Gamma s$，以 $s(0) = 1$ 积分便是上述指数律。
振幅阻尼信道同时把失去的激发布居转移给基态，保持总概率为一。
有限环境发生复现时，未发射比例一般不满足这个乘法关系。

\begin{theorem}{有限维GKSL生成元}{qm-gksl}
    有限维系统的连续、时间齐次完全正迹保持半群，其生成元可写为
    \begin{equation}
        \mathcal L(\op{\rho}) = -\frac\i\hbar[\op{H}_\eff,\op{\rho}]
        + \sum_\mu\Bigp{\op{L}_\mu\op{\rho} \op{L}_\mu\adjoint
            - \frac12\{\op{L}_\mu\adjoint \op{L}_\mu,\op{\rho}\}}.
        \label{eq:lindblad}
    \end{equation}
    其中$\op{H}_\eff$自伴，$\op{L}_\mu$是系统算符。
    反过来，这一形式生成完全正迹保持半群。
\end{theorem}
GKSL 是 Gorini、Kossakowski、Sudarshan 与 Lindblad 的姓氏缩写。
本书将速率平方根吸收到$\op{L}_\mu$中，其量纲为时间$^{-1/2}$。
这一有限维分类定理的证明见文献\cite{manzano}。

从短时 Kraus 算符可看清反对易项的来历。先将各通道的贡献合写成算符
\[
    \op{C} = \sum_\mu \op{L}_\mu\adjoint \op{L}_\mu .
\]
保留无跃迁算符的一阶项和跃迁算符的最低阶项，取
\begin{subequations}
    \begin{align}
        \op{K}_0 &= \one - \biggp{\frac{\i \op{H}_\eff}\hbar + \frac{\op{C}}{2}}\d t,\\
        \op{K}_\mu &= \sqrt{\d t}\,\op{L}_\mu.
    \end{align}
\end{subequations}
代入Kraus和并保留$\d t$，便得到式\eqref{eq:lindblad}。

分别展开无跃迁项和跃迁项，得到
\begin{subequations}
    \begin{alignat}{1}
        \op{K}_0\op{\rho} \op{K}_0\adjoint &= \op{\rho} - \frac{\i}{\hbar}[\op{H}_\eff,\op{\rho}]\d t
            - \frac12(\op{C}\op{\rho} + \op{\rho} \op{C})\d t + \bigo(\d t^2),\\
        \sum_\mu \op{K}_\mu\op{\rho} \op{K}_\mu\adjoint &= \sum_\mu \op{L}_\mu\op{\rho} \op{L}_\mu\adjoint\d t.
    \end{alignat}
\end{subequations}
相加后减去$\op{\rho}$、除以$\d t$，得到主方程。
其中$\op{H}_\eff$为自伴的有效Hamilton算符，
可包含环境引起的能移；耗散由$\op{L}_\mu$表示。

概率守恒要求所有记录的概率之和为一。由上面的短时算符可验算
\[
    \op{K}_0\adjoint \op{K}_0 + \sum_\mu \op{K}_\mu\adjoint \op{K}_\mu
    = \one - \op{C}\d t + \op{C}\d t + \bigo(\d t^2) = \one + \bigo(\d t^2).
\]
无跃迁项中 $\op{C}$ 的系数恰好使一阶概率变化相消。
具体的$\op{H}_\eff$和$\op{L}_\mu$由系统与环境的耦合及以下近似确定。

微观推导通常先假定系统与环境弱耦合，采用 Born 近似。再要求环境关联时间远短于系统弛豫(relaxation)时间，采用 Markov 近似。
随后对可分辨的不同跃迁频率作长期平均，即久期近似(secular approximation)；也可以使用其他保证完全正性的受控极限。
这些是不同步骤。强耦合、结构化谱和长期记忆可能导致非Markov动力学；
未经久期处理的Redfield方程也不自动保持正性\cite{manzano}。

\section{布居弛豫与退相干}
同一条主方程怎样给出寿命实验和 Ramsey 实验各自测到的衰减时间？
原子发射一个光子会减少激发态的数目，能级间相位的随机变化则可以在布居不变时使条纹变淡。
把这两种过程分别代入上节的主方程，便能分别跟踪激发态的概率和条纹的对比度。

设 $\op{H} = \hbar\omega_0\ketbra{\excitedstate}{\excitedstate}$，$\ket{\groundstate}$ 为基态。零温自发发射和纯退相干(pure dephasing)分别取跳跃算符
\begin{equation}
    \op{L}_1 = \sqrt{\Gamma}\ketbra{\groundstate}{\excitedstate},\qquad
    \op{L}_\upphi = \sqrt{\frac{\gamma_\upphi}{2}}\,\op{\sigma}_z .
\end{equation}
其中 $\Gamma$ 是自发发射率，$\gamma_\upphi$ 是纯退相干率。

$\rho_{\excitedstate\excitedstate} = \mel{\excitedstate}{\op{\rho}}{\excitedstate}$是激发态布居，
$\rho_{\excitedstate\groundstate} = \mel{\excitedstate}{\op{\rho}}{\groundstate}$是基态与激发态之间的相干，
$\op{\sigma}_z = \ketbra{\excitedstate}{\excitedstate} - \ketbra{\groundstate}{\groundstate}$。
两种耗散的算符形式分别为
\begin{subequations}
    \begin{alignat}{1}
        \mathcal D_1(\op{\rho}) &= \Gamma\rho_{\excitedstate\excitedstate}\ketbra{\groundstate}{\groundstate}
            - \frac{\Gamma}{2}\{\ketbra{\excitedstate}{\excitedstate},\op{\rho}\},\\
        \mathcal D_\upphi(\op{\rho}) &= \frac{\gamma_\upphi}{2}(\op{\sigma}_z\op{\rho}\op{\sigma}_z - \op{\rho}).
    \end{alignat}
\end{subequations}
第一式的馈入项只增加$\rho_{\groundstate\groundstate}$，
反对易项在$\excitedstate\excitedstate$元贡献$-\Gamma\rho_{\excitedstate\excitedstate}$，
在$\excitedstate\groundstate$元贡献$-\Gamma\rho_{\excitedstate\groundstate}/2$。
第二式对对角元为零，对$\excitedstate\groundstate$元贡献$-\gamma_\upphi\rho_{\excitedstate\groundstate}$。

直接代入矩阵元得
\begin{subequations}
    \begin{alignat}{1}
        \dot\rho_{\excitedstate\excitedstate} &= -\Gamma\rho_{\excitedstate\excitedstate},\\
        \dot\rho_{\excitedstate\groundstate} &= -\Bigp{\i\omega_0 + \frac\Gamma2 + \gamma_\upphi}\rho_{\excitedstate\groundstate}.
    \end{alignat}
\end{subequations}
\begin{definition}{布居弛豫时间与相干时间}{qm-relaxation-times}
    若能量布居与平衡值之差按指数衰减，其幅值降至初值 $1/\e$ 的时间称为 $T_1$。
    若能量基中非对角元的模按指数衰减，其幅值降至初值 $1/\e$ 的时间称为 $T_2$。
    在理想 Ramsey 读出中，后一衰减可由条纹对比度测得。
\end{definition}
读出上述两个微分方程的衰减率，得到这个零温两态模型中的关系
\begin{subequations}
    \begin{gather}
        T_1 = \Gamma^{-1},\qquad
        T_2^{-1} = \frac1{2T_1} + \gamma_\upphi,\\
        T_2 \leq 2T_1.
    \end{gather}
\end{subequations}
纯退相干打乱两能级之间的相对相位而保持布居，因此只对 $T_2^{-1}$ 贡献附加的一项（图\ref{fig:relaxation}）。
\begin{figure}[htbp]
    \centering
    \includegraphics[width=.62\linewidth]{figures/relaxation.pdf}
    \caption{各量分别除以自身初值，时间以 $T_1$ 为单位。虚线为激发布居 $\rho_{\excitedstate\excitedstate}$；两条实线为相干幅值 $\abs{\rho_{\excitedstate\groundstate}}$。取 $\gamma_\upphi T_1 = 0,1$ 时，相应的衰减率比 $T_1/T_2$ 分别为 $0.5,1.5$。}
    \label{fig:relaxation}
\end{figure}

解这两个一阶微分方程可得
\begin{subequations}
    \begin{alignat}{1}
        \rho_{\excitedstate\excitedstate}(t) &= \rho_{\excitedstate\excitedstate}(0)\e^{-t/T_1},\\
        \rho_{\excitedstate\groundstate}(t) &= \rho_{\excitedstate\groundstate}(0)\e^{-\i\omega_0t-t/T_2},\\
        \rho_{\groundstate\groundstate}(t) &= 1 - \rho_{\excitedstate\excitedstate}(t).
    \end{alignat}
\end{subequations}
第三式由迹守恒确定，显示激发态损失的布居进入了基态。

有热激发时需再加$\sqrt{\Gamma_{\up}}\ket{\excitedstate}\bra{\groundstate}$，
布居弛豫率变为$\Gamma_\downarrow + \Gamma_{\up}$，平衡布居由两率之比决定。

Ramsey实验在理想脉冲下读出
\begin{equation}
    P_+(T) = \frac12[1 + \e^{-T/T_2}\cos(\Delta T + \phi_0)]
\end{equation}
$\phi_0$是两次脉冲设置的相对相位；
这里取初始制备与读出轴使纵向布居项不贡献该信号。
静态失谐分布的系综平均还会引入$T_2^*$；其包络未必指数，且可被自旋回波部分恢复。
回波可用来区分静态失谐引起的相位散布与不可逆的相干衰减。

\begin{example}{静态失谐的自由感应与自旋回波}{qm-spin-echo}
    设每个粒子的附加失谐$\delta$在一次实验内固定，
    不同粒子的$\delta$服从均值为零、标准差为$\sigma_\delta$的 Gauss 分布。
    求自由等待 $T$ 后的系综相干因子，再比较在 $T/2$ 时绕 $x$ 轴加理想 $\pi$ 脉冲的情形。

    \par\smallskip\textbf{解}\quad 自由等待 $T$ 后的系综相干含因子
    \begin{equation}
        \intR
        \frac{\e^{-\delta^2/(2\sigma_\delta^2)}}{\sqrt{2\pi}\sigma_\delta}
        \e^{-\i\delta T}\d\delta
        = \e^{-\sigma_\delta^2T^2/2}.
    \end{equation}
    Gauss 包络的宽度由静态失谐的标准差 $\sigma_\delta$ 决定。
    若在$T/2$插入理想$\pi$脉冲$\op{R}_x(\pi) = -\i\op{\sigma}_x$，
    利用$\op{\sigma}_x\op{\sigma}_z\op{\sigma}_x = -\op{\sigma}_z$，
    两段失谐演化满足
    \begin{equation}
        \e^{-\i\delta T\op{\sigma}_z/4}
        \op{R}_x(\pi)
        \e^{-\i\delta T\op{\sigma}_z/4}
        = \op{R}_x(\pi).
    \end{equation}
    固定失谐的相位因而被消去，末端适当读出可恢复对比度。
    在两半段之间变化的噪声及真实能量弛豫通常不能这样完全消除。
\end{example}

\section{量子跳跃}
上节的布居曲线平均了许多次制备。若在一次实验中也记下光子到达的时刻，
每次探测到光子都提供了一次发射事件的时间信息。如何根据这类记录更新原子态，并由单次记录恢复平均衰减曲线？

\begin{definition}{光子计数下的量子跳跃}{qm-photon-jump}
    在理想光子计数中，某一发射通道的点击所引起的条件态突变称为量子跳跃(quantum jump)。
    对以 $\op{L}_1$ 表示的通道，短时间 $\d t$ 内的点击概率（取一阶）和点击后的态分别为
    \begin{subequations}
        \begin{align}
            p_1 &= \d t\,\tr(\op{L}_1\adjoint \op{L}_1\op{\rho}),\\
            \op{\rho}_1 &= \frac{\op{L}_1\op{\rho} \op{L}_1\adjoint}{\tr(\op{L}_1\adjoint \op{L}_1\op{\rho})}\quad(p_1 > 0).
        \end{align}
    \end{subequations}
\end{definition}
若所有发射通道均受到理想监测且无其他耗散，没有点击时按 $\op{K}_0$ 更新并归一化。
将各记录的条件态按概率加权，便恢复主方程。所选的监测方式决定单次记录对应的条件轨迹。
探测效率不足时，“没有计数”还可能表示已经发射、但光子漏检。
计算这组记录的条件态时，需要将未发射与漏检两种可能按概率合并。

\sectionstar{短时演化与量子Zeno效应}
光子计数等待发射自行发生。若改为周期性测量系统是否仍处于初态，每次测量的结果还决定了下一时段的初态。
缩短测量间隔会怎样影响生存概率，取决于单次演化中离开初态的概率如何随时间增长。
这里要回到封闭系统的幺正演化，考察指数衰减模型尚未分辨的时间尺度。

令 $\op{H}$ 为不显含时间的自伴 Hamilton 算符。对归一化初态 $\ket{\Psi}$，生存振幅(survival amplitude) $A(t)$ 是演化后的态在原态方向上的分量。它是复数，由下列矩阵元给出，
\begin{equation}\label{eq:qm-survival-amplitude}
    A(t) = \mel{\Psi}{\exp(-\i \op{H}t/\hbar)}{\Psi}.
\end{equation}
\begin{theorem}{短时生存律}{qm-short-time-survival}
    封闭系统按不显含时间的自伴 $\op{H}$ 演化，若归一化初态 $\ket{\Psi}$ 的能量方差有限，
    则在时刻 $t$ 测得系统仍处于该初态的概率（生存概率，survival probability）满足
    \begin{equation}
        \abs{\mel{\Psi}{\e^{-\i \op{H}t/\hbar}}{\Psi}}^2
        = 1 - \frac{(\Delta H)^2t^2}{\hbar^2} + o(t^2).
    \end{equation}
\end{theorem}
\begin{proof}
    能量方差有限时，能量分布的二阶矩存在。对 $A(t)$ 的能量谱积分求出在 $t = 0$ 的前两阶导数，便有
    \begin{equation}
        A(t) = 1 - \frac{\i\ave{\op{H}}t}{\hbar}
        - \frac{\ave{\op{H}^2}t^2}{2\hbar^2} + o(t^2).
    \end{equation}
    与其共轭相乘，一阶项相消，二阶系数为 $-(\ave{\op{H}^2} - \ave{\op{H}}^2)/\hbar^2$。
\end{proof}
因此在极短时间内损失的概率以 $t^2$ 起始。指数律 $\e^{-\Gamma t} = 1 - \Gamma t + \cdots$ 则描述环境相关时间之后的有效动力学。
\begin{corollary}[理想量子 Zeno 效应]
在上述条件下，于固定总时间 $t$ 内作 $N$ 次等间隔的理想瞬时测量，每次都投影到同一个一维初态及其正交补。
每次回答“仍在初态”的分支都从 $\ket{\Psi}$ 重新出发，连续 $N$ 次都得到这一结果的概率为
\begin{equation}
    P_N(t) = \abs{A(t/N)}^{2N}
        = \Bigb{1 - \frac{(\Delta H)^2t^2}{N^2\hbar^2} + o(N^{-2})}^N \to 1.
\end{equation}
\end{corollary}
实际能达到的生存概率还取决于每次测量的持续时间、带宽与反馈。

\begin{example}{原子衰变中的概率守恒}{qm-decay-trace}
    激发态布居随时间衰减。比较全系综密度算符的迹与只保留尚未发射原子的未归一化密度算符的迹。

    \par\smallskip\textbf{解}\quad 激发态失去的布居由跳跃项馈入基态，两者的迹变化相消，总迹保持为一。若只保留尚未发射的子系综，其未归一化态的迹给出该记录的概率。
\end{example}
\reading{完全正映射和GKSL形式核对Manzano的教学综述\cite{manzano}。
    最后一章回到精密光谱，说明非相对论单粒子理论何时必须扩展。}
